% 教学构造：一个图块4个灰度箱、8个像素，计数上限2；右图为四个图块中心与插值点(alpha=.25,beta=.5)。
% 三个面板并排，总宽约 108mm。
\begin{tikzpicture}[x=4mm,y=3.2mm,font=\figurefont\small]
  % 面板一：截断
  \foreach \k in {0,...,3}{\node[below,font=\figurefont\footnotesize] at (\k,0) {\k};}
  \fill[FigureBlueFill,draw=FigureBlue,line width=.6pt] (.62,0) rectangle (1.38,2);
  \fill[FigureRedFill,draw=FigureRed,line width=.6pt] (.62,2) rectangle (1.38,8);
  \draw[FigureStroke,line width=.6pt] (-.6,0)--(3.6,0);
  \draw[FigureYellow,dashed,line width=.9pt] (-.6,2)--(3.6,2);
  \node[anchor=west,font=\figurefont\footnotesize] at (3.6,2) {上限};
  \node[anchor=west,font=\figurefont\footnotesize,text=FigureMuted] at (1.5,6) {截下6};
  \node at (1.5,-2.4) {截断前};
  % 面板二：重新分配
  \begin{scope}[xshift=30mm]
    \foreach \k in {0,...,3}{
      \fill[FigureBlueFill,draw=FigureBlue,line width=.6pt] (\k-.38,0) rectangle (\k+.38,2);
      \node[below,font=\figurefont\footnotesize] at (\k,0) {\k};
    }
    \draw[FigureStroke,line width=.6pt] (-.6,0)--(3.6,0);
    \draw[FigureYellow,dashed,line width=.9pt] (-.6,2)--(3.6,2);
    \node at (1.5,-2.4) {重新分配后};
  \end{scope}
  % 面板三：图块与双线性插值
  \begin{scope}[xshift=64mm,x=1mm,y=1mm]
    \draw[FigureGrid,line width=.5pt] (0,0) grid[step=12] (24,24);
    \foreach \cx/\cy/\v in {6/18/120,18/18/140,6/6/100,18/6/160}{
      \fill[FigureBlue] (\cx,\cy) circle (1pt);
      \node[font=\figurefont\footnotesize,anchor=south] at (\cx,\cy+.8) {\v};
    }
    \draw[FigureMuted,dashed,line width=.5pt] (6,18) rectangle (18,6);
    \fill[FigureRed] (9,12) circle (1.2pt);
    \node at (12,-4.6) {四个图块的映射};
  \end{scope}
\end{tikzpicture}
